\documentclass[11pt]{article}
\usepackage{amsmath,amssymb,amsthm}
\usepackage[margin=1in]{geometry}

\newtheorem{theorem}{Theorem}
\newtheorem{lemma}[theorem]{Lemma}

\title{Problem 435: Polynomials of Fibonacci Numbers}
\author{Project Euler Solutions}
\date{}

\begin{document}
\maketitle

\section{Problem Statement}
Let $F_n$ be the $n$-th Fibonacci number ($F_0 = 0$, $F_1 = 1$). Define $P_k(x) = \sum_{n=0}^{k} F_n x^n$. Evaluate $\sum_{n=1}^{N} P_k(n) \pmod{p}$ for given $k$, $N$, and prime $p$.

\section{Mathematical Foundation}

\begin{theorem}[Binet's Formula]
For $n \geq 0$,
\[
F_n = \frac{\varphi^n - \psi^n}{\sqrt{5}},
\]
where $\varphi = \frac{1+\sqrt{5}}{2}$ and $\psi = \frac{1-\sqrt{5}}{2}$.
\end{theorem}

\begin{proof}
The characteristic equation $x^2 - x - 1 = 0$ has roots $\varphi, \psi$. The general solution $F_n = A\varphi^n + B\psi^n$ with initial conditions $F_0 = 0$, $F_1 = 1$ gives $A = 1/\sqrt{5}$, $B = -1/\sqrt{5}$.
\end{proof}

\begin{lemma}[Closed Form for $P_k(x)$]
For $\varphi x \neq 1$ and $\psi x \neq 1$:
\[
P_k(x) = \frac{1}{\sqrt{5}} \left( \frac{(\varphi x)^{k+1} - 1}{\varphi x - 1} - \frac{(\psi x)^{k+1} - 1}{\psi x - 1} \right).
\]
\end{lemma}

\begin{proof}
Substitute Binet's formula and evaluate each geometric sum $\sum_{n=0}^{k} r^n = (r^{k+1}-1)/(r-1)$.
\end{proof}

\begin{theorem}[Existence of $\sqrt{5}$ in $\mathbb{F}_p$]
For an odd prime $p \neq 5$, $\sqrt{5}$ exists in $\mathbb{F}_p$ iff $p \equiv \pm 1 \pmod{5}$.
\end{theorem}

\begin{proof}
By quadratic reciprocity, $\left(\frac{5}{p}\right) = \left(\frac{p}{5}\right)$. The quadratic residues modulo $5$ are $\{1, 4\}$, so $5$ is a QR mod $p$ iff $p \equiv 1$ or $4 \pmod{5}$.
\end{proof}

\begin{lemma}[Tonelli--Shanks]
Given a quadratic residue $a$ modulo an odd prime $p$, the Tonelli--Shanks algorithm computes $r$ with $r^2 \equiv a \pmod{p}$ in $O(\log^2 p)$ time.
\end{lemma}

\begin{proof}
Factor $p - 1 = 2^s q$ with $q$ odd. The algorithm iteratively refines a candidate root using the $2$-Sylow subgroup structure of $(\mathbb{Z}/p\mathbb{Z})^*$, converging in at most $s \leq \log_2 p$ iterations.
\end{proof}

\begin{theorem}[Summation Formula]
\[
\sum_{n=1}^{N} P_k(n) \equiv \frac{1}{\sqrt{5}} \sum_{n=1}^{N} \left( \frac{(\varphi n)^{k+1} - 1}{\varphi n - 1} - \frac{(\psi n)^{k+1} - 1}{\psi n - 1} \right) \pmod{p}.
\]
Each term is computed in $O(\log k + \log p)$ via modular exponentiation and inverse.
\end{theorem}

\begin{proof}
Direct substitution of Lemma~1 into the sum. Modular inverses exist since $\varphi n \not\equiv 1 \pmod{p}$ for all but $O(1)$ values of $n$.
\end{proof}

\section{Algorithm}
\begin{enumerate}
    \item Compute $\sqrt{5} \bmod p$ via Tonelli--Shanks; derive $\varphi, \psi \bmod p$.
    \item For each $n = 1, \ldots, N$: evaluate $P_k(n) \bmod p$ using the closed form.
    \item Accumulate the sum modulo $p$.
\end{enumerate}

\section{Complexity Analysis}
\begin{itemize}
    \item \textbf{Time:} $O(N \log k)$ per-term exponentiation; reducible to $O(\log N \cdot \log k)$ via matrix methods.
    \item \textbf{Space:} $O(1)$ auxiliary.
\end{itemize}

\section{Answer}

\[
\boxed{252541322550}
\]

\end{document}
